\vspace*{\credgap}
{\noindent\LARGE\sc Letter from the Editor}
\vspace{\credgap}

\vspace{\ackgap}\noindent
Dear Reader,

\vspace{\credgap}\noindent
Before you sits the newest issue of Aporia, the Feminist Edition. After several months of deliberation, it is with great pride that I present the following papers to you. They are both a testament to their writers and the Femporia team who have worked with great diligence to bring together this issue. All of this would not be possible without the amenability of the writers, the potential of the works themselves, and you, the reader, to whom we now pass the torch.

We live in reductive times; wherein more and more processes are being shorted and made facile by artificial intelligence. Motivation to engage with any texts in their natural state, unprocessed, is actively discouraged by the ease of generative summaries and critical synopses. Even when it comes to our own essays, we are likely to consult some kind of computer to offer some words of encouragement. This in turn means that away dwindles passion in undergraduate students for essay writing, especially aside from that required of them.

In light of this technology fuelled benightedness, it is now more important than ever to keep reading and writing, no matter how small the act may seem. We may just be a university philosophy journal but still, student led publications like this one are an important access point for writer and reader alike. They provide us undergraduates with feedback, editing and publishing opportunities.

And in whatever contribution you may have given, as a reader, writer or editor, you have a right to be proud. We are all responsible for the preservation of our own discussion and engagement.

With that being said, it is with gratitude and sorrow that I present my first and last issue of The Feminist Edition. I am incredibly thankful for this opportunity, and I look forward to seeing what the next generation of editors gets up to. Although my time as editor-in-chief may have been short-lived, as the proverb goes, I hope that this issue still burns bright. 

\vspace{\credgap}\noindent
That’s all from me, 

\vspace{\ackgap}\noindent
Ella Johnston  

\noindent
Editor-in-Chief, \emph{Aporia: The Feminist Edition}
